\documentclass[12pt,a4paper]{article}
\usepackage{cmap} % Makes the PDF copiable. See http://tex.stackexchange.com/a/64198/25761
\usepackage[T1]{fontenc}
\usepackage[brazil]{babel}
\usepackage[utf8]{inputenc}
\usepackage{amsmath}
\usepackage{amsfonts}
\usepackage{amssymb}
\usepackage{amsthm}
\usepackage{textcomp} % \degree
\usepackage{gensymb} % \degree
\usepackage[usenames,svgnames,dvipsnames]{xcolor}
\usepackage{hyperref}
\usepackage{graphicx}
\usepackage[margin=2cm]{geometry}

\hypersetup{
    colorlinks = true,
    allcolors = {blue}
}

\newcommand*\sen{\operatorname{sen}}
\newcommand*\senh{\operatorname{senh}}

\newcommand*\tipo{PROVA II}
\newcommand*\disciplina{CDI0001}
\newcommand*\eu{Helder G. G. de Lima}
\newcommand*\data{16/10/2015}

\author{\eu}
\title{\tipo - \disciplina}
\date{\data}

\begin{document}
\thispagestyle{empty}
\newgeometry{margin=2cm,bottom=0.5cm}
\begin{center}
\includegraphics[width=9.0cm]{marca}
\noindent\begin{tabular}{l c c r}
  \textbf{\disciplina (turma: \rule{4cm}{0.01cm})}
& \textbf{\tipo}
& \textbf{\data}
\end{tabular}
\\ Prof. \eu\footnote{
Este é um material de acesso livre distribuído sob os termos da licença \href{https://creativecommons.org/licenses/by-sa/4.0/deed.pt_BR}{Creative Commons BY-SA 4.0}.}
\end{center}

\noindent Nome do(a) aluno(a): \rule{13cm}{0.01cm}

%\section*{Instruções}
\begin{center}\fbox{
\begin{minipage}{14cm}

{\footnotesize
\begin{itemize}
\renewcommand{\theenumi}{\Roman{enumi}}
\item Identifique-se em todas as folhas.
\item Mantenha o celular e os demais equipamentos eletrônicos desligados durante a prova.
\item Justifique cada resposta com cálculos ou argumentos baseados na teoria estudada.
\item Resolva (integralmente) apenas os itens de que precisar para somar 10,0 pontos.
\end{itemize}
}

\end{minipage}
}
\end{center}


%\section*{Questões}

\begin{enumerate}
\item (0,5) Dê exemplos de funções $f, g: \mathbb{R} \to \mathbb{R}$ tais que $\left[f( g(x) )\right]^\prime \neq f^\prime(x)g^\prime(x)$. Faça as contas para mostrar que realmente são diferentes.

\item (1,0) Seja $f(x) = \dfrac{10^{2x}}{\ln(10)}$. Determine uma fórmula para $f^{(n)}$, a derivada de ordem $n$ de $f$.

\item (0,5) Calcule a derivada de $y = \senh(x)\cosh(x) + \sen(\cos(\ln(\pi + 1)))$.

\item (2,0) Uma escada de $250cm$ de comprimento está apoiada em um muro vertical. Se a extremidade inferior da escada se afasta do muro na razão de $3cm/s$, então com que rapidez está descendo a extremidade superior no instante em que o pé da escada está a $70cm$ do muro?

\item (1,5) Mostre que se
$y = f(x) = e^x \sen(x)$
então vale
$y^{\prime \prime} - 2y^\prime +2y = 0$,
sendo $y^\prime$ e $y^{\prime\prime}$, respectivamente, as derivadas de primeira e segunda ordem da função indicada.

\item (1,5) Utilize uma aproximação linear para estimar o valor de $\sqrt{26}$.

\item (2,5) Seja $y=f(x)$ dada implicitamente pela equação $e^{(y^2 \cdot x)}=x-y$. Calcule $f^\prime(0)$.

\item (2,5) Determine a equação da reta normal ao gráfico de $f(x) = \tan(x)$ em $x=\pi/4$.

\end{enumerate}

\newpage
\restoregeometry
\section*{Respostas e observações}

\begin{enumerate}
\item \textit{ Dê exemplos de funções $f, g: \mathbb{R} \to \mathbb{R}$ tais que $\left[f( g(x) )\right]^\prime \neq f^\prime(x)g^\prime(x)$. Faça as contas para mostrar que realmente são diferentes.}


Há inúmeras funções para as quais isso ocorre. Por exemplo:
\begin{itemize}
\item Se $f(x) = x^2$ e $g(x) = 3x$, então:
\[
\left[f(g(x))\right]^\prime
= \left[(3x)^2\right]^\prime
= \left[9x^2\right]^\prime
= 18x
\neq 6x
= 2x \cdot 3
= (x^2)^\prime \cdot (3x)^\prime
= f^\prime(x)g^\prime(x).
\]
\item Se $f(x) = \sen(x)$ e $g(x) = e^x$, então:
\[
\left[f(g(x))\right]^\prime
= \left[\sen(e^x)\right]^\prime
= \cos(e^x)e^x
\neq \cos(x) \cdot e^x
= (\sen(x))^\prime \cdot (e^x)^\prime
= f^\prime(x)g^\prime(x).
\]
\end{itemize}

\item \textit{ Seja $f(x) = \dfrac{10^{2x}}{\ln(10)}$. Determine uma fórmula para $f^{(n)}$, a derivada de ordem $n$ de $f$.}

Como $\left[ 10^{2x} \right]^\prime = \left[ 10^{2x} \cdot \ln(10) \cdot 2 \right] = 2\ln(10) \cdot 10^{2x}$, tem-se
\[
f^\prime(x)
= \dfrac{ 2\ln(10) \cdot 10^{2x} }{\ln(10)}
= 2 \cdot 10^{2x}
\]
\[
f^{(2)}(x)
= 2 \cdot \left[ 10^{2x} \right]^\prime
= 2 \cdot \left[ 2\ln(10) \cdot 10^{2x} \right]
= 2^2 \ln(10) \cdot 10^{2x}
\]
\[
f^{(3)}(x)
= 2^2 \ln(10) \cdot \left[ 10^{2x} \right]^\prime
= 2^2 \ln(10) \cdot \left[ 2\ln(10) \cdot 10^{2x} \right]
= 2^3 \ln(10)^2 \cdot 10^{2x}
\]

Pode-se concluir que a derivada de ordem $n$ será $f^{(n)}(x)= 2^n \ln(10)^{n-1} \cdot 10^{2x}$.


\item \textit{ Calcule a derivada de $y = \senh(x)\cosh(x) + \sen(\cos(\ln(\pi + 1)))$. }

Como $\sen(\cos(\ln(\pi + 1)))$ é apenas uma constante, a sua derivada é zero. Então:
\[
y^\prime
= (\senh(x)\cosh(x))^\prime + 0
= (\senh(x))^\prime\cosh(x) + \senh(x) (\cosh(x))^\prime
%= \cosh(x) \cosh(x) + \senh(x) \senh(x)
= \cosh^2(x) + \senh^2(x).
\]

\item \textit{ Uma escada de $250cm$ de comprimento está apoiada em um muro vertical. Se a extremidade inferior da escada se afasta do muro na razão de $3cm/s$, então com que rapidez está descendo a extremidade superior no instante em que o pé da escada está a $70cm$ do muro? }

Se $x$ é a distância (em centímetros) entre o muro e o pé da escada, e $y$ a distância (em centímetros) entre o chão e o topo da escada, então $x^2+y^2=250^2$. Derivando ambos os membros da igualdade em relação a $t$ obtem-se $2x \cdot x^\prime + 2y \cdot y^\prime = 0$, que equivale a $y^\prime = -\frac{2x}{2y} \cdot x^\prime = -\frac{x}{y} \cdot x^\prime$.

Quando $x = 70cm$, tem-se $70^2+y^2=250^2$ e portanto $y=\sqrt{ 250^2 - 70^2 } = 10 \sqrt{ 25^2 - 7^2} = 10 \cdot 24 = 240 cm$. Consequentemente, $y^\prime = -\frac{70}{240} \cdot 3 = -7/8 = -0,875 cm/s$.

\item \textit{ Mostre que se
$y = f(x) = e^x \sen(x)$
então vale
$y^{\prime \prime} - 2y^\prime +2y = 0$,
sendo $y^\prime$ e $y^{\prime\prime}$, respectivamente, as derivadas de primeira e segunda ordem da função indicada. }

Tem-se
\[
y^\prime
= (e^x \sen(x))^\prime
= e^x \sen(x) + e^x \cos(x)
= e^x \left[ \sen(x) + \cos(x) \right]
\]
\[
y^{\prime\prime}
= (e^x \left[ \sen(x) + \cos(x) \right] )^\prime
= e^x\left[ \sen(x) + \cos(x) \right] + e^x ( \cos(x) - \sen(x) )
= 2 e^x \cos(x)
\]
Consequentemente,
\[
y^{\prime \prime} - 2y^\prime +2y
= 2 e^x \cos(x)
- 2 \left( e^x \left[ \sen(x) + \cos(x) \right] \right)
+ 2 \left( e^x \sen(x) \right)
= 0.
\]

\item \textit{ Utilize uma aproximação linear para estimar o valor de $\sqrt{26}$. }


Seja $f(x) = \sqrt{x}$, $x_0 = 25$ e $\Delta x = 1$. Então e $f^\prime = \dfrac{1}{2\sqrt{x}}$ e tem-se:

\begin{align*}
\sqrt{26} = \sqrt{25 + 1}
= f(x_0+\Delta x)
& \approx f(x_0) + f^{\prime}(x_0)\Delta x \\
& = \sqrt{25} + \frac{1}{2\sqrt{25}} \cdot 1
= 5 + \frac{1}{10}
= 5 + 0,1
= 5,1.
\end{align*}


\item \textit{ Seja $y=f(x)$ dada implicitamente pela equação $e^{(y^2 \cdot x)}=x-y$. Calcule $f^\prime(0)$. }

Derivando ambos os membros da equação em relação a $x$, resulta que
\[
e^{(y^2 \cdot x)}(2y y^\prime x + y^2)=1-y^\prime.
\]
Em particular, se $x=0$ tem-se $e^{(y^2 \cdot 0)}(2y y^\prime \cdot 0 + y^2)=1-y^\prime$, isto é, $y^2=1-y^\prime$. Além disso, o valor de $y$ correspondente pode ser obtido levando em conta que $e^{(y^2 \cdot 0)}=0-y$, ou seja, $y=-1$. Assim, $(-1)^2=1-y^\prime$ implica que $y^\prime = 0$.

\item \textit{ Determine a equação da reta normal ao gráfico de $f(x) = \tan(x)$ em $x=\pi/4$. }

Tendo em mente que $\tan(\pi/4) = 1$, que $f^\prime(x) = (\tan(x))^\prime = \sec^2(x)$, e que
\[
\sec^2(\pi/4)
= \frac{1}{\cos^2(\pi/4)}
= \frac{1}{ \left(\frac{ \sqrt{2} }{2}\right)^2 }
= \frac{1}{ \left(\frac{ 2 }{4}\right) }
= 2,
\]
conclui-se que a equação da reta tangente é $y - 1 = 2(x-\pi/4)$, ou seja, $y = 2x + (1-\pi/2)$.
\end{enumerate}
\end{document}
